
\documentclass{article}
\usepackage{colortbl}
\usepackage{makecell}
\usepackage{multirow}
\usepackage{supertabular}

\begin{document}

\newcounter{utterance}

\centering \large Interaction Transcript for game `eqbench', experiment `eq\_bench', episode 106 with reflexion{-}Qwen3.5{-}4B.
\vspace{24pt}

{ \footnotesize  \setcounter{utterance}{1}
\setlength{\tabcolsep}{0pt}
\begin{supertabular}{c@{$\;$}|p{.15\linewidth}@{}p{.15\linewidth}p{.15\linewidth}p{.15\linewidth}p{.15\linewidth}p{.15\linewidth}}
    \# & \multicolumn{2}{c}{Player} && \multicolumn{2}{c}{Game Master} \\
    \hline

    \theutterance \stepcounter{utterance}  
    & & & \multicolumn{4}{p{0.6\linewidth}}{
        \cellcolor[rgb]{0.9,0.9,0.9}{
            \makecell[{{p{\linewidth}}}]{
                \texttt{\tiny{[P1$\langle$GM]}}
                \texttt{You have attempted this same task before and did not succeed. These are the plans you wrote after those attempts.} \\
\texttt{Trial \#1: New plan: I misinterpreted the emotional state of Terry based on the dialogue. Terry is a protective, guilt{-}ridden character who just threatened to bench Cass to keep her alive. While he feels frustrated with Cass's recklessness, the act of making a hard decision to protect her life suggests a release of tension and a sense of duty fulfilled, which aligns more with "Relieved" than "Frustrated." I should have recognized that the threat "If it means keeping you alive, yes" indicates a resolution to a conflict, leading to relief rather than lingering frustration. I will re{-}evaluate the final emotional beat as a moment of relief for Terry after making a difficult but necessary choice.} \\
\\ 
\texttt{Frustrated: 5} \\
\texttt{Relieved: 9} \\
\texttt{Dismissive: 2} \\
\texttt{Content: 0} \\
\\ 
\\ 
\\ 
\texttt{Your task is to predict the likely emotional responses of a character in this dialogue:} \\
\\ 
\texttt{Terry: We're not going in blind, Cass. Not this time.} \\
\texttt{Cass: And who's gonna stop me? You?} \\
\texttt{Terry: If I have to. } \\
\texttt{Cass: Oh, right. The big, bad Terry Holloway. Always with a plan. } \\
\texttt{Terry: You're goddamn right. You know why? Because plans save lives.} \\
\texttt{Cass: Or maybe they just make you feel better about yourself. Like you're in control.} \\
\texttt{Terry: This isn't about control, Cass. It's about survival. } \\
\texttt{Cass: Survival? Or just your guilt?} \\
\texttt{Terry: That's enough.} \\
\texttt{Cass: Hit a nerve, did I? } \\
\texttt{Terry: You're good, Cass. But you're reckless. And that's a liability I can't afford.} \\
\texttt{Cass: And there it is. The threat. You gonna bench me, Terry?} \\
\texttt{Terry: If it means keeping you alive, yes.} \\
\texttt{Cass: I don't need your protection. I need to do my job.} \\
\texttt{{[}End dialogue{]}} \\
\\ 
\texttt{At the end of this dialogue, Terry would feel...} \\
\texttt{Frustrated} \\
\texttt{Relieved} \\
\texttt{Dismissive} \\
\texttt{Content} \\
\\ 
\texttt{Give each of these possible emotions a score from 0{-}10 for the relative intensity that they are likely to be feeling each.} \\
\\ 
\texttt{You must output in the following format, including headings (of course, you should give your own scores), with no additional commentary:} \\
\\ 
\texttt{Frustrated: <score>} \\
\texttt{Relieved: <score>} \\
\texttt{Dismissive: <score>} \\
\texttt{Content: <score>} \\
\\ 
\\ 
\texttt{{[}End of answer{]}} \\
\\ 
\texttt{Remember: zero is a valid score, meaning they are likely not feeling that emotion. You must score at least one emotion > 0.} \\
\\ 
\texttt{Your answer:} \\
            }
        }
    }
     \\ \\

    \theutterance \stepcounter{utterance}  
    & \multicolumn{4}{p{0.6\linewidth}}{
        \cellcolor[rgb]{0.561,0.737,0.561}{
            \makecell[{{p{\linewidth}}}]{
                \texttt{\tiny{[P1$\rangle$GM]}}
                \texttt{Frustrated: 8} \\
\texttt{Relieved: 3} \\
\texttt{Dismissive: 2} \\
\texttt{Content: 0} \\
            }
        }
    }
    & & \\ \\

    \theutterance \stepcounter{utterance}  
    & & & \multicolumn{2}{p{0.3\linewidth}}{
        \cellcolor[rgb]{0.9,0.9,0.9}{
            \makecell[{{p{\linewidth}}}]{
                \texttt{\tiny{[GM$|$GM]}}
                \texttt{Frustrated: 8} \\
\texttt{Relieved: 3} \\
\texttt{Dismissive: 2} \\
\texttt{Content: 0} \\
            }
        }
    }
    & & \\ \\

    \theutterance \stepcounter{utterance}  
    & & & \multicolumn{2}{p{0.3\linewidth}}{
        \cellcolor[rgb]{0.9,0.9,0.9}{
            \makecell[{{p{\linewidth}}}]{
                \texttt{\tiny{[GM$|$GM]}}
                \texttt{Frustrated: 8} \\
\texttt{Relieved: 0} \\
\texttt{Dismissive: 5} \\
\texttt{Content: 0} \\
            }
        }
    }
    & & \\ \\

    \theutterance \stepcounter{utterance}  
    & & & \multicolumn{2}{p{0.3\linewidth}}{
        \cellcolor[rgb]{0.9,0.9,0.9}{
            \makecell[{{p{\linewidth}}}]{
                \texttt{\tiny{[GM$|$GM]}}
                \texttt{game\_result = LOSE} \\
            }
        }
    }
    & & \\ \\

    \theutterance \stepcounter{utterance}  
    & & \multicolumn{4}{p{0.6\linewidth}}{
        \cellcolor[rgb]{0.95,0.95,0.95}{
            \makecell[{{p{\linewidth}}}]{
                \texttt{\tiny{[P1$\leftrightarrow$P]}}
                \texttt{Post{-}episode Reflexion (not sent to GM):} \\
\texttt{New plan: The previous attempt correctly identified "Relieved" as the primary emotion but underestimated the intensity of "Frustrated." Terry's dialogue reveals a deep{-}seated frustration with Cass's recklessness ("You're reckless. And that's a liability I can't afford") which is the direct cause of his decision to bench her. While the act of making the final decision brings relief, the underlying emotion driving the entire interaction is frustration with her behavior. I will adjust the scores to reflect that the frustration is high (driving the action) while the relief is moderate (the result of the action). I will also ensure "Dismissive" remains low as Terry is engaging deeply with Cass's arguments, not dismissing them.} \\
\\ 
\texttt{Frustrated: 9} \\
\texttt{Relieved: 6} \\
\texttt{Dismissive: 2} \\
\texttt{Content: 0} \\
            }
        }
    }
    & \\ \\

\end{supertabular}
}

\end{document}
